\chapter{Notation and Conventions}

\begin{remarkbox}
This appendix collects the notational conventions used throughout the notes. It
is not a substitute for the main text. Its purpose is to make signs, indices,
Fourier conventions, differential-form notation, and gamma-matrix conventions
explicit in one place.
\end{remarkbox}

\section{Metric Signatures}

Throughout these notes, unless stated otherwise, flat Minkowski spacetime is
written with the mostly-minus convention
\[
    \eta_{\mu\nu}=\operatorname{diag}(+1,-1,-1,-1).
\]
Thus the invariant interval is
\[
    ds^2=\eta_{\mu\nu}dx^\mu dx^\nu=dt^2-d\vect{x}^{\,2}.
\]
The square of a four-vector \(p^\mu=(E,\vect{p})\) is
\[
    p_\mu p^\mu=E^2-|\vect{p}|^2.
\]
For a massive relativistic excitation,
\[
    p_\mu p^\mu=m^2.
\]

Partial derivatives are denoted by
\[
    \partial_\mu=\frac{\partial}{\partial x^\mu}.
\]
With the mostly-minus metric,
\[
    \partial^0=\partial_0=\partial_t,
    \qquad
    \partial^i=-\partial_i.
\]
The d'Alembert operator is
\[
    \Box=\partial_\mu\partial^\mu=\partial_t^2-\nabla^2.
\]

\begin{definitionbox}
The convention used in these notes is
\[
    \eta_{\mu\nu}=\operatorname{diag}(+---),
    \qquad
    \Box=\partial_t^2-\nabla^2.
\]
Changing signature changes several signs in field equations and actions.
\end{definitionbox}

Curved spacetime metrics are usually denoted by \(g_{\mu\nu}\). In that case,
indices are raised and lowered with \(g^{\mu\nu}\) and \(g_{\mu\nu}\), and the
volume density is \(\sqrt{|g|}\,d^{d+1}x\), where \(g=\det(g_{\mu\nu})\).

\section{Index Conventions}

Greek indices usually denote spacetime indices:
\[
    \mu,\nu,\rho,\sigma=0,1,2,3.
\]
Latin indices from the middle of the alphabet usually denote spatial indices:
\[
    i,j,k=1,2,3.
\]
Internal indices are often written as
\[
    a,b,c=1,\ldots,\dim\mathfrak{g}
\]
for a Lie algebra \(\mathfrak{g}\). Field-component labels may be written as
\(A,B,C\).

The Einstein summation convention is used: repeated upper-lower index pairs are
summed over. For example,
\[
    A_\mu B^\mu=\sum_{\mu=0}^3 A_\mu B^\mu.
\]
Dummy indices may be renamed. Free indices must match on both sides of an
equation. For instance,
\[
    \partial_\mu F^{\mu\nu}=J^\nu
\]
is an equation with one free index \(\nu\). The index \(\mu\) is summed.

For Euclidean spatial vectors, bold notation is used:
\[
    \vect{x},\qquad \vect{E},\qquad \vect{B}.
\]
The Euclidean dot product is written
\[
    \vect{a}\cdot\vect{b}=a_i b_i.
\]

\section{Units and Dimensions}

Most field-theory formulas are written in natural units,
\[
    c=\hbar=1.
\]
In these units, time and length have the same dimension, and mass, energy,
frequency, and inverse length have the same dimension. For example,
\[
    [m]=[E]=[\omega]=[L^{-1}].
\]

The action is dimensionless in units \(\hbar=1\):
\[
    [S]=1.
\]
Therefore the Lagrangian density in \(D=d+1\) spacetime dimensions has mass
dimension
\[
    [\mathcal{L}]=D.
\]
For a real scalar field with kinetic term
\[
    \frac{1}{2}\partial_\mu\phi\partial^\mu\phi,
\]
one obtains
\[
    [\phi]=\frac{D-2}{2}.
\]
In four spacetime dimensions, \([\phi]=1\).

Dimensional analysis is not a decorative check. It is often the fastest way to
find missing constants, wrong powers, or inconsistent interaction terms.

\section{Fourier Transform Conventions}

For spatial Fourier transforms, a common convention used in these notes is
\[
    f(\vect{x})=\int\frac{d^d k}{(2\pi)^d}\,\widetilde{f}(\vect{k})
    e^{i\vect{k}\cdot\vect{x}},
\]
with inverse
\[
    \widetilde{f}(\vect{k})=\int d^d x\,f(\vect{x})e^{-i\vect{k}\cdot\vect{x}}.
\]
For spacetime plane waves, we often write
\[
    e^{-ik_\mu x^\mu}=e^{-i\omega t+i\vect{k}\cdot\vect{x}}.
\]
With the mostly-minus convention,
\[
    k_\mu x^\mu=\omega t-\vect{k}\cdot\vect{x}.
\]

For the Klein--Gordon equation,
\[
    (\Box+m^2)\phi=0,
\]
plane waves obey
\[
    \omega^2=|\vect{k}|^2+m^2.
\]

\begin{remarkbox}
Fourier conventions vary from book to book. The placement of signs in the
exponential and factors of \(2\pi\) must always be checked before comparing
propagators or Green's functions across references.
\end{remarkbox}

\section{Levi-Civita Symbols and Tensors}

The three-dimensional Levi-Civita symbol is denoted by \(\epsilon_{ijk}\), with
\[
    \epsilon_{123}=+1.
\]
It is antisymmetric in all indices. The cross product is
\[
    (\vect{a}\times\vect{b})_i=\epsilon_{ijk}a_jb_k.
\]

In four dimensions, the Levi-Civita symbol is chosen as
\[
    \epsilon^{0123}=+1
\]
unless stated otherwise. Care is required when lowering indices in Lorentzian
signature. With the mostly-minus metric, signs may change when all indices are
lowered.

The dual electromagnetic field strength is
\[
    \widetilde{F}^{\mu\nu}=\frac{1}{2}\epsilon^{\mu\nu\rho\sigma}F_{\rho\sigma}.
\]
Useful Lorentz invariants are
\[
    F_{\mu\nu}F^{\mu\nu}=2(|\vect{B}|^2-|\vect{E}|^2),
\]
and, with the conventions of these notes,
\[
    F_{\mu\nu}\widetilde{F}^{\mu\nu}=-4\vect{E}\cdot\vect{B}.
\]

\section{Differential Form Conventions}

A \(p\)-form is written as
\[
    \alpha=\frac{1}{p!}\alpha_{\mu_1\cdots\mu_p}
    dx^{\mu_1}\wedge\cdots\wedge dx^{\mu_p}.
\]
The wedge product is antisymmetric:
\[
    dx^\mu\wedge dx^\nu=-dx^\nu\wedge dx^\mu.
\]
The exterior derivative satisfies
\[
    d^2=0.
\]

For electromagnetism,
\[
    A=A_\mu dx^\mu,
    \qquad
    F=dA.
\]
For non-Abelian gauge theory,
\[
    F=dA+gA\wedge A.
\]
The Abelian Bianchi identity is
\[
    dF=0,
\]
while the Yang--Mills Bianchi identity is
\[
    DF=0.
\]

\section{Gamma Matrix Conventions}

Gamma matrices satisfy
\[
    \{\gamma^\mu,\gamma^\nu\}=2\eta^{\mu\nu}\mathbb{I}.
\]
With the mostly-minus convention,
\[
    (\gamma^0)^2=\mathbb{I},
    \qquad
    (\gamma^i)^2=-\mathbb{I}.
\]
The Dirac adjoint is
\[
    \bar\psi=\psi^\dagger\gamma^0.
\]
The free Dirac equation is
\[
    (i\gamma^\mu\partial_\mu-m)\psi=0.
\]
The Dirac current is
\[
    j^\mu=\bar\psi\gamma^\mu\psi.
\]

In four spacetime dimensions,
\[
    \gamma^5=i\gamma^0\gamma^1\gamma^2\gamma^3,
\]
and the chiral projectors are
\[
    P_L=\frac{1}{2}(\mathbb{I}-\gamma^5),
    \qquad
    P_R=\frac{1}{2}(\mathbb{I}+\gamma^5).
\]

\begin{summarybox}
The global convention of these notes is mostly-minus signature, natural units,
Einstein summation, \(\Box=\partial_t^2-\nabla^2\), and gamma matrices satisfying
\(\{\gamma^\mu,\gamma^\nu\}=2\eta^{\mu\nu}\mathbb{I}\). Whenever comparing with a
different reference, check metric signature, Fourier signs, gauge-coupling
conventions, and gamma-matrix conventions first.
\end{summarybox}
